\section{Motivation}

Teams of autonomous robots are poised to transform domains ranging from disaster
response to planetary exploration. A swarm of small, inexpensive robots can search
a collapsed building faster than a single sophisticated platform; a fleet of autonomous
underwater vehicles can map an ocean basin with greater coverage and resilience
than any individual sensor package. Yet realizing these benefits depends on solving
a fundamental challenge: how should a collection of independent agents coordinate
their actions when each agent has only a partial, noisy view of the world and
cannot reliably communicate with its teammates?

Centralized approaches --- where a single planner computes actions for all agents ---
offer theoretical simplicity and, in some cases, optimality guarantees. However,
they suffer from three well-known weaknesses: they introduce a single point of
failure, they scale poorly with team size, and they assume a communication
infrastructure that is unavailable in many practical settings. The alternative,
decentralized coordination, distributes decision-making across agents and promises
fault tolerance, scalability, and communication efficiency. The cost is a more
difficult planning problem, because each agent must reason not only about the
environment but also about the beliefs, intentions, and likely future actions of
its teammates.

This dissertation addresses that cost. We develop a framework that makes
decentralized coordination both principled and practical for teams of robots
operating in uncertain, dynamic environments.
